% Figure 01 — Architecture globale
% Internet -> Route53 -> EIP -> EC2 -> Docker stack
\documentclass[border=10pt]{standalone}
% Préambule commun pour tous les diagrammes TikZ standalone
% À inclure via % Préambule commun pour tous les diagrammes TikZ standalone
% À inclure via % Préambule commun pour tous les diagrammes TikZ standalone
% À inclure via \input{preambule.tex} après \documentclass{standalone}
\usepackage[utf8]{inputenc}
\usepackage[T1]{fontenc}
\usepackage[french]{babel}
\usepackage{lmodern}
\usepackage{tikz}
\usepackage{xcolor}

\usetikzlibrary{
  arrows.meta,
  positioning,
  shapes.geometric,
  shapes.symbols,
  shapes.misc,
  fit,
  calc,
  backgrounds,
  decorations.pathreplacing,
  decorations.markings,
  shadows.blur,
  patterns
}

% Palette de couleurs (cohérente sur tout le rapport)
\definecolor{awsorange}{HTML}{FF9900}
\definecolor{awsdark}{HTML}{232F3E}
\definecolor{awsblue}{HTML}{1A476F}
\definecolor{dockerblue}{HTML}{2496ED}
\definecolor{dockerdark}{HTML}{0DB7ED}
\definecolor{nginxgreen}{HTML}{009639}
\definecolor{flaskgray}{HTML}{3776AB}
\definecolor{pgblue}{HTML}{336791}
\definecolor{tlsgreen}{HTML}{2E8B57}
\definecolor{netgray}{HTML}{6C757D}
\definecolor{bgsoft}{HTML}{F5F7FA}
\definecolor{bordersoft}{HTML}{C8D1DC}
\definecolor{textdark}{HTML}{1F2937}
\definecolor{accentred}{HTML}{C0392B}
\definecolor{accentyellow}{HTML}{F1C40F}

% Styles génériques réutilisés dans plusieurs diagrammes
\tikzset{
  base/.style={
    rounded corners=4pt,
    draw=bordersoft,
    fill=white,
    text=textdark,
    align=center,
    inner sep=6pt,
    font=\small
  },
  zone/.style={
    base,
    fill=bgsoft,
    draw=netgray,
    dashed,
    inner sep=12pt
  },
  cloud/.style={
    base,
    fill=awsdark!5,
    draw=awsorange,
    line width=0.8pt,
    inner sep=14pt
  },
  service/.style={
    base,
    line width=0.6pt,
    minimum width=2.6cm,
    minimum height=1.1cm,
    drop shadow={shadow xshift=1pt, shadow yshift=-1pt, opacity=0.2}
  },
  legendebox/.style={
    base,
    fill=white,
    font=\footnotesize,
    inner sep=4pt
  },
  fleche/.style={
    -{Latex[length=2.5mm,width=2mm]},
    line width=0.6pt,
    color=netgray
  },
  fleche-bidir/.style={
    {Latex[length=2.5mm,width=2mm]}-{Latex[length=2.5mm,width=2mm]},
    line width=0.6pt,
    color=netgray
  },
  fleche-bold/.style={
    -{Latex[length=3mm,width=2.4mm]},
    line width=1pt
  },
  etiquette/.style={
    font=\scriptsize\sffamily,
    fill=white,
    inner sep=2pt,
    text=textdark
  }
}
 après \documentclass{standalone}
\usepackage[utf8]{inputenc}
\usepackage[T1]{fontenc}
\usepackage[french]{babel}
\usepackage{lmodern}
\usepackage{tikz}
\usepackage{xcolor}

\usetikzlibrary{
  arrows.meta,
  positioning,
  shapes.geometric,
  shapes.symbols,
  shapes.misc,
  fit,
  calc,
  backgrounds,
  decorations.pathreplacing,
  decorations.markings,
  shadows.blur,
  patterns
}

% Palette de couleurs (cohérente sur tout le rapport)
\definecolor{awsorange}{HTML}{FF9900}
\definecolor{awsdark}{HTML}{232F3E}
\definecolor{awsblue}{HTML}{1A476F}
\definecolor{dockerblue}{HTML}{2496ED}
\definecolor{dockerdark}{HTML}{0DB7ED}
\definecolor{nginxgreen}{HTML}{009639}
\definecolor{flaskgray}{HTML}{3776AB}
\definecolor{pgblue}{HTML}{336791}
\definecolor{tlsgreen}{HTML}{2E8B57}
\definecolor{netgray}{HTML}{6C757D}
\definecolor{bgsoft}{HTML}{F5F7FA}
\definecolor{bordersoft}{HTML}{C8D1DC}
\definecolor{textdark}{HTML}{1F2937}
\definecolor{accentred}{HTML}{C0392B}
\definecolor{accentyellow}{HTML}{F1C40F}

% Styles génériques réutilisés dans plusieurs diagrammes
\tikzset{
  base/.style={
    rounded corners=4pt,
    draw=bordersoft,
    fill=white,
    text=textdark,
    align=center,
    inner sep=6pt,
    font=\small
  },
  zone/.style={
    base,
    fill=bgsoft,
    draw=netgray,
    dashed,
    inner sep=12pt
  },
  cloud/.style={
    base,
    fill=awsdark!5,
    draw=awsorange,
    line width=0.8pt,
    inner sep=14pt
  },
  service/.style={
    base,
    line width=0.6pt,
    minimum width=2.6cm,
    minimum height=1.1cm,
    drop shadow={shadow xshift=1pt, shadow yshift=-1pt, opacity=0.2}
  },
  legendebox/.style={
    base,
    fill=white,
    font=\footnotesize,
    inner sep=4pt
  },
  fleche/.style={
    -{Latex[length=2.5mm,width=2mm]},
    line width=0.6pt,
    color=netgray
  },
  fleche-bidir/.style={
    {Latex[length=2.5mm,width=2mm]}-{Latex[length=2.5mm,width=2mm]},
    line width=0.6pt,
    color=netgray
  },
  fleche-bold/.style={
    -{Latex[length=3mm,width=2.4mm]},
    line width=1pt
  },
  etiquette/.style={
    font=\scriptsize\sffamily,
    fill=white,
    inner sep=2pt,
    text=textdark
  }
}
 après \documentclass{standalone}
\usepackage[utf8]{inputenc}
\usepackage[T1]{fontenc}
\usepackage[french]{babel}
\usepackage{lmodern}
\usepackage{tikz}
\usepackage{xcolor}

\usetikzlibrary{
  arrows.meta,
  positioning,
  shapes.geometric,
  shapes.symbols,
  shapes.misc,
  fit,
  calc,
  backgrounds,
  decorations.pathreplacing,
  decorations.markings,
  shadows.blur,
  patterns
}

% Palette de couleurs (cohérente sur tout le rapport)
\definecolor{awsorange}{HTML}{FF9900}
\definecolor{awsdark}{HTML}{232F3E}
\definecolor{awsblue}{HTML}{1A476F}
\definecolor{dockerblue}{HTML}{2496ED}
\definecolor{dockerdark}{HTML}{0DB7ED}
\definecolor{nginxgreen}{HTML}{009639}
\definecolor{flaskgray}{HTML}{3776AB}
\definecolor{pgblue}{HTML}{336791}
\definecolor{tlsgreen}{HTML}{2E8B57}
\definecolor{netgray}{HTML}{6C757D}
\definecolor{bgsoft}{HTML}{F5F7FA}
\definecolor{bordersoft}{HTML}{C8D1DC}
\definecolor{textdark}{HTML}{1F2937}
\definecolor{accentred}{HTML}{C0392B}
\definecolor{accentyellow}{HTML}{F1C40F}

% Styles génériques réutilisés dans plusieurs diagrammes
\tikzset{
  base/.style={
    rounded corners=4pt,
    draw=bordersoft,
    fill=white,
    text=textdark,
    align=center,
    inner sep=6pt,
    font=\small
  },
  zone/.style={
    base,
    fill=bgsoft,
    draw=netgray,
    dashed,
    inner sep=12pt
  },
  cloud/.style={
    base,
    fill=awsdark!5,
    draw=awsorange,
    line width=0.8pt,
    inner sep=14pt
  },
  service/.style={
    base,
    line width=0.6pt,
    minimum width=2.6cm,
    minimum height=1.1cm,
    drop shadow={shadow xshift=1pt, shadow yshift=-1pt, opacity=0.2}
  },
  legendebox/.style={
    base,
    fill=white,
    font=\footnotesize,
    inner sep=4pt
  },
  fleche/.style={
    -{Latex[length=2.5mm,width=2mm]},
    line width=0.6pt,
    color=netgray
  },
  fleche-bidir/.style={
    {Latex[length=2.5mm,width=2mm]}-{Latex[length=2.5mm,width=2mm]},
    line width=0.6pt,
    color=netgray
  },
  fleche-bold/.style={
    -{Latex[length=3mm,width=2.4mm]},
    line width=1pt
  },
  etiquette/.style={
    font=\scriptsize\sffamily,
    fill=white,
    inner sep=2pt,
    text=textdark
  }
}


\begin{document}
\begin{tikzpicture}[node distance=8mm and 14mm]

  % --- Bloc Internet / Client ---
  \node[service, fill=netgray!10, draw=netgray, minimum width=3.2cm]
    (client) {\textbf{Navigateur Web}\\\footnotesize Utilisateur final};

  \node[service, fill=accentyellow!15, draw=accentyellow!80!black,
        right=18mm of client, minimum width=3cm]
    (dns) {\textbf{Route\,53}\\\footnotesize \texttt{webcyber.app}};

  % --- Cloud AWS ---
  \node[cloud, right=22mm of dns,
        minimum width=8.5cm, minimum height=7.5cm,
        label={[font=\small\bfseries, text=awsorange]north:AWS Cloud — eu-west-3 (Paris)}]
    (aws) {};

  % VPC à l'intérieur du cloud AWS
  \node[zone, fill=awsblue!5, draw=awsblue,
        minimum width=7cm, minimum height=6cm,
        label={[font=\footnotesize\bfseries, text=awsblue]north:VPC \texttt{10.0.0.0/16}}]
    at ([yshift=-3mm]aws.center) (vpc) {};

  % Security Group
  \node[zone, fill=accentred!5, draw=accentred,
        minimum width=5.6cm, minimum height=4.6cm,
        label={[font=\footnotesize\bfseries, text=accentred]north:Security Group \texttt{webcyber-sg}}]
    at ([yshift=-4mm]vpc.center) (sg) {};

  % EC2
  \node[service, fill=awsorange!15, draw=awsorange,
        minimum width=4.6cm, minimum height=3.4cm,
        font=\small,
        label={[font=\scriptsize\bfseries, text=awsdark]north:EC2 \texttt{t2.micro} — Ubuntu 22.04}]
    at ([yshift=-3mm]sg.center) (ec2) {};

  % Docker stack à l'intérieur de l'EC2
  \node[base, fill=dockerblue!10, draw=dockerblue,
        minimum width=4cm, minimum height=2.4cm,
        font=\scriptsize,
        label={[font=\scriptsize\bfseries, text=dockerblue]north:Docker Compose}]
    at ([yshift=-2mm]ec2.center) (docker) {};

  % 3 conteneurs miniatures
  \node[base, fill=nginxgreen!20, draw=nginxgreen, font=\tiny,
        minimum width=1cm, minimum height=0.7cm]
    at ([xshift=-12mm,yshift=-2mm]docker.center) (nginx) {Nginx};
  \node[base, fill=flaskgray!25, draw=flaskgray, font=\tiny,
        minimum width=1cm, minimum height=0.7cm]
    at ([xshift=0mm,yshift=-2mm]docker.center) (app) {Flask};
  \node[base, fill=pgblue!25, draw=pgblue, font=\tiny,
        minimum width=1cm, minimum height=0.7cm]
    at ([xshift=12mm,yshift=-2mm]docker.center) (db) {Postgres};

  % Elastic IP en étiquette sur EC2
  \node[etiquette, draw=awsorange, fill=awsorange!15,
        anchor=south east]
    at ([xshift=-2mm,yshift=2mm]ec2.south east) {Elastic IP};

  % --- Flèches ---
  \draw[fleche-bold, color=netgray]
    (client) -- node[etiquette, above]{1. DNS} (dns);

  \draw[fleche-bold, color=tlsgreen]
    (dns.east) to[bend left=8]
    node[etiquette, above, text=tlsgreen]{2. HTTPS \texttt{:443}}
    (sg.west);

  \draw[fleche-bold, color=tlsgreen, dashed]
    (client.south) to[bend right=22]
    node[etiquette, below, text=tlsgreen]{réponse chiffrée TLS}
    (sg.south west);

  % Légende des ports SG
  \node[legendebox, anchor=north west, font=\tiny,
        align=left, fill=accentred!5, draw=accentred]
    at ([xshift=2mm,yshift=-2mm]sg.north west)
    {\textbf{Ports ouverts}\\
     22/tcp \, 80/tcp \, 443/tcp\\
     source : \texttt{0.0.0.0/0}};

\end{tikzpicture}
\end{document}
